\usepackage{float,csquotes,textcomp}
\usepackage[backend=bibtex,style=authoryear,maxcitenames=2]{biblatex}
\addbibresource{referencias.bib}
\let\cleardoublepage\clearpage
\setcounter{tocdepth}{1}
\setcounter{secnumdepth}{1}
\setlength{\headheight}{28pt}
\fancyhead[L]{\small\textcolor{UEAPDarkGreen}{Oficina de análise ecológica | UEAP}}
\fancyhead[R]{\small\textcolor{UEAPDarkGreen}{Caderno de aula}}
\fancypagestyle{plain}{\fancyhf{}\fancyfoot[L]{\small Robson Borges de Lima}\fancyfoot[R]{\small\thepage}\renewcommand{\headrulewidth}{0pt}}
\titleformat{\chapter}[display]{\normalfont\bfseries\color{UEAPDarkGreen}}
 {\Large PARTE \thechapter}{.25cm}{\fontsize{24}{28}\selectfont}[\vspace{.15cm}\titlerule]
\titlespacing*{\chapter}{0pt}{12pt}{22pt}
\lstdefinestyle{aulaR}{style=ueapRstyle,
 basicstyle=\ttfamily\fontsize{9}{11}\selectfont,
 numbers=left,numberstyle=\ttfamily\scriptsize\color{UEAPGray},
 numbersep=9pt,xleftmargin=2.6em,xrightmargin=3pt,
 frame=l,rulecolor=\color{UEAPDarkGreen},framesep=7pt,
 breaklines=true,breakatwhitespace=false,columns=fullflexible,upquote=true,
 aboveskip=8pt,belowskip=8pt,
 literate={→}{{$\rightarrow$}}1}
\newcommand{\codigo}[2]{\begingroup\singlespacing\lstinputlisting[style=aulaR,
 firstline=#1,lastline=#2,firstnumber=#1]{../atividade_docente.R}\endgroup}
\newcommand{\pergunta}[1]{\begin{exercicios}#1\end{exercicios}}
\newcommand{\prancha}[4]{
 \ifnum\value{section}>0\clearpage\fi\section{#1}\label{sec:#2}
 \begin{figure}[H]\centering
 \includegraphics[width=\linewidth,height=.59\textheight,keepaspectratio]{../saidas/figuras/#2.pdf}
 \caption{#3}\label{fig:#2}\end{figure}
 \begin{conceito}[title={Leitura para discussão}]#4\end{conceito}}
\hypersetup{pdftitle={ANOVA, teste F, Tukey e regressão em ecologia: caderno de aula},
 pdfauthor={Robson Borges de Lima},linkcolor=UEAPDarkBlue}
\autor{Robson Borges de Lima}
\disciplina{Estatística Aplicada}
\newcommand{\capaatividade}{
\hypersetup{pageanchor=false}
\begin{titlepage}\thispagestyle{empty}
\begin{tikzpicture}[remember picture,overlay]
\fill[UEAPDarkGreen](current page.north west)rectangle([xshift=1.05cm]current page.south west);
\fill[UEAPGold]([xshift=1.05cm]current page.north west)rectangle([xshift=1.15cm]current page.south west);
\node[rotate=90,text=white,font=\bfseries\scriptsize]at([xshift=.52cm]current page.west){UEAP\quad|\quad CADERNO DE AULA};
\end{tikzpicture}
\hspace*{.6cm}\begin{minipage}{.94\textwidth}
\begin{minipage}{.72\linewidth}
{\Large\bfseries\color{UEAPDarkGreen}Universidade do Estado do Amapá}\\[8pt]
{\small Colegiado de Engenharia Florestal\\Programa de Pós-graduação em\\Recursos Naturais Amazônicos}
\end{minipage}\hfill\begin{minipage}{.22\linewidth}\includegraphics[width=\linewidth]{Brasao_UEAP.png}\end{minipage}
\vspace{1.8cm}\par
{\small\bfseries\color{UEAPGold}ESTATÍSTICA APLICADA | ATIVIDADE PRÁTICA}\par\vspace{.5cm}
{\fontsize{28}{33}\selectfont\bfseries\color{UEAPDarkBlue}ANOVA, teste F,\\Tukey e regressão\\em ecologia\par}
\vspace{.65cm}{\large\color{UEAPDarkGreen}Exploração, inferência e recursos gráficos em R}\par
\vspace{.6cm}{\color{UEAPGold}\rule{.85\linewidth}{1.2pt}}\par
\vspace{.5cm}{\large\bfseries Prof. Dr. Robson Borges de Lima}\par
\vspace{.3cm}{\small Aula de 27 de outubro de 2026 | RENAmazon}\par
\vspace{1cm}
\begin{tcolorbox}[colback=UEAPLightGray,colframe=UEAPDarkGreen,boxrule=.5pt,arc=1mm]
\textbf{Caderno para acompanhamento em aula.} Roteiro de discussão, fundamentos, resultados interpretados, 19 pranchas gráficas e script integral com numeração contínua de linhas.
\end{tcolorbox}
\vspace{.6cm}{\small Macapá, Amapá | 2026\\Dados simulados para ensino | Versão 1.0}
\end{minipage}\end{titlepage}\hypersetup{pageanchor=true}}
